% GENERATED by `python -m loanshield.report` from report/results_template.tex.
% Do not edit by hand: edit the template and re-run.
% =====================================================================
% LoanShield - Results section (Springer report), TEMPLATE.
%
% Do not edit the generated results_section.tex by hand. Edit this file, then run
%     python -m loanshield.report
% to fill the double-angle-bracket placeholders from results/ and data/ and copy
% the figures into report/figures/.
%
% Markers in the generated file:
%   \MISSING{key}  a number not produced yet (run the experiments)
%   \AUTHOR{...}   interpretation the team must write once numbers are final
% Preamble needs: \usepackage{booktabs,graphicx,xcolor,multirow}
% =====================================================================
\providecommand{\MISSING}[1]{\textcolor{red}{\texttt{[#1]}}}
\providecommand{\AUTHOR}[1]{\textcolor{red}{[AUTHOR: #1]}}

\section{Results}\label{sec:results}

All results use developer-grouped five-fold cross-validation unless stated otherwise, so apps from one developer never appear on both sides of a split. The positive class is \emph{unverified}. We report PR-AUC (average precision) as the primary metric because the classes are imbalanced, together with ROC-AUC, recall at a 5\% false-positive rate, and precision, recall and F1 for the unverified class (threshold chosen on the training folds). Values are mean $\pm$ standard deviation over folds.

Headline numbers are computed on \emph{organic} apps, i.e.\ apps discovered through Play Store search or similar-app links. Apps found through the Play links in the RBI directory itself are verified by construction, so including them would reward a model for recognising how an app was found rather than its risk. All models are trained on every labelled app; results on all apps are reported alongside for comparison.

% ---------------------------------------------------------------------
\subsection{Dataset and labels}\label{sec:res-data}

The RBI Digital Lending Apps directory snapshot (2026-10-08) contained 2,869 entries from 457 regulated entities, linking 434 distinct Play Store apps. Play Store discovery produced 1,149 candidate apps (345 from search, 526 from similar-app links, 278 from RBI-listed links), of which 627 were Finance apps that mention lending. Table~\ref{tab:data} summarises the labelled dataset.

\begin{table}[t]
\centering
\caption{Labelled dataset. \emph{Organic} apps were found by Play search or similar-app links.}\label{tab:data}
\begin{tabular}{lrr}
\toprule
 & All apps & Organic apps \\
\midrule
Verified (listed in the RBI DLA directory) & 351 & 163 \\
Unverified (lends, not listed) & 27 & 27 \\
Base rate (share unverified) & 0.071 & 0.142 \\
Developers & 341 & 169 \\
\midrule
Excluded (does not lend) & 119 & \\
Unique reviews (labelled apps) & 44,064 & \\
\bottomrule
\end{tabular}
\end{table}

Automatic matching against the RBI directory labelled 378 apps: 341 verified by an exact Play package match, 10 verified by app-name and developer match, and 27 unverified because they offer credit but match no directory entry. A further 0 came from the flagged-app seed list and 0 from manual review. Two annotators independently labelled the 130-app review queue, with Cohen's $\kappa = $ \MISSING{data.kappa} (\MISSING{data.agreement} raw agreement). Unverified apps are markedly newer: 15 of the 27 were released in 2026, consistent with Google Play's removal of unlisted loan apps from January 2026.

% ---------------------------------------------------------------------
\subsection{Model comparison}\label{sec:res-main}

Table~\ref{tab:main} compares the tabular, text and fusion models on organic apps. The best tabular model was \MISSING{main.best\_tabular} and the best text model \MISSING{main.best\_text}; the fusion model stacks the two.

\begin{table}[t]
\centering
\caption{Developer-grouped 5-fold CV on organic apps (base rate 0.142). The last column gives PR-AUC on all apps (base rate 0.071).}\label{tab:main}
\setlength{\tabcolsep}{4pt}
\begin{tabular}{llcccccc}
\toprule
 & Model & PR-AUC & ROC-AUC & Rec@5\%FPR & Precision & Recall & PR-AUC (all) \\
\midrule
\multirow{4}{*}{Tabular}
 & Logistic regression & \MISSING{main.organic.logreg.pr\_auc} & \MISSING{main.organic.logreg.roc\_auc} & \MISSING{main.organic.logreg.recall\_at\_5fpr} & \MISSING{main.organic.logreg.precision} & \MISSING{main.organic.logreg.recall} & \MISSING{main.all.logreg.pr\_auc} \\
 & Random forest & \MISSING{main.organic.rf.pr\_auc} & \MISSING{main.organic.rf.roc\_auc} & \MISSING{main.organic.rf.recall\_at\_5fpr} & \MISSING{main.organic.rf.precision} & \MISSING{main.organic.rf.recall} & \MISSING{main.all.rf.pr\_auc} \\
 & XGBoost & \MISSING{main.organic.xgb.pr\_auc} & \MISSING{main.organic.xgb.roc\_auc} & \MISSING{main.organic.xgb.recall\_at\_5fpr} & \MISSING{main.organic.xgb.precision} & \MISSING{main.organic.xgb.recall} & \MISSING{main.all.xgb.pr\_auc} \\
 & LightGBM & \MISSING{main.organic.lgbm.pr\_auc} & \MISSING{main.organic.lgbm.roc\_auc} & \MISSING{main.organic.lgbm.recall\_at\_5fpr} & \MISSING{main.organic.lgbm.precision} & \MISSING{main.organic.lgbm.recall} & \MISSING{main.all.lgbm.pr\_auc} \\
\midrule
\multirow{3}{*}{Text}
 & TF-IDF + LR & \MISSING{main.organic.tfidf\_lr.pr\_auc} & \MISSING{main.organic.tfidf\_lr.roc\_auc} & \MISSING{main.organic.tfidf\_lr.recall\_at\_5fpr} & \MISSING{main.organic.tfidf\_lr.precision} & \MISSING{main.organic.tfidf\_lr.recall} & \MISSING{main.all.tfidf\_lr.pr\_auc} \\
 & TF-IDF + SVM & \MISSING{main.organic.tfidf\_svm.pr\_auc} & \MISSING{main.organic.tfidf\_svm.roc\_auc} & \MISSING{main.organic.tfidf\_svm.recall\_at\_5fpr} & \MISSING{main.organic.tfidf\_svm.precision} & \MISSING{main.organic.tfidf\_svm.recall} & \MISSING{main.all.tfidf\_svm.pr\_auc} \\
 & MuRIL & \MISSING{main.organic.muril.pr\_auc} & \MISSING{main.organic.muril.roc\_auc} & \MISSING{main.organic.muril.recall\_at\_5fpr} & \MISSING{main.organic.muril.precision} & \MISSING{main.organic.muril.recall} & \MISSING{main.all.muril.pr\_auc} \\
\midrule
Fusion & \MISSING{main.fusion\_name} & \MISSING{main.organic.fusion.pr\_auc} & \MISSING{main.organic.fusion.roc\_auc} & \MISSING{main.organic.fusion.recall\_at\_5fpr} & \MISSING{main.organic.fusion.precision} & \MISSING{main.organic.fusion.recall} & \MISSING{main.all.fusion.pr\_auc} \\
\bottomrule
\end{tabular}
\end{table}

\AUTHOR{Compare the best model with the logistic-regression baseline and with the base rate; say whether tree models, text models or fusion help, and whether the fold standard deviations make the differences meaningful.}

\begin{figure}[t]
\centering
\includegraphics[width=0.48\linewidth]{figures/pr_curves_grouped.png}\hfill
\includegraphics[width=0.48\linewidth]{figures/confusion_grouped.png}
\caption{Left: precision--recall curves on organic apps (out-of-fold scores; dashed line = base rate). Right: confusion matrices summed over folds.}\label{fig:pr}
\end{figure}

% ---------------------------------------------------------------------
\subsection{Generalisation to newer apps}\label{sec:res-temporal}

To test whether models trained on older apps flag newly released ones, we trained on apps released before \MISSING{temporal.cutoff} (\MISSING{temporal.n\_train} apps, \MISSING{temporal.n\_train\_pos} unverified) and tested on later releases (\MISSING{temporal.n\_test} organic apps, \MISSING{temporal.n\_test\_pos} unverified; base rate \MISSING{temporal.base\_rate}). Table~\ref{tab:temporal} reports the three best models.

\begin{table}[t]
\centering
\caption{Temporal hold-out on organic apps released on or after \MISSING{temporal.cutoff}.}\label{tab:temporal}
\begin{tabular}{lccc}
\toprule
Model & PR-AUC & ROC-AUC & Rec@5\%FPR \\
\midrule
\MISSING{temporal.1.model} & \MISSING{temporal.1.pr\_auc} & \MISSING{temporal.1.roc\_auc} & \MISSING{temporal.1.recall\_at\_5fpr} \\
\MISSING{temporal.2.model} & \MISSING{temporal.2.pr\_auc} & \MISSING{temporal.2.roc\_auc} & \MISSING{temporal.2.recall\_at\_5fpr} \\
\MISSING{temporal.3.model} & \MISSING{temporal.3.pr\_auc} & \MISSING{temporal.3.roc\_auc} & \MISSING{temporal.3.recall\_at\_5fpr} \\
\bottomrule
\end{tabular}
\end{table}

\AUTHOR{Compare with the grouped-CV results. If performance drops, relate it to the change in the unverified population after the January 2026 Play removal (older and newer unverified apps differ).}

% ---------------------------------------------------------------------
\subsection{Feature-group ablation}\label{sec:res-ablation}

Removing one feature group at a time from \MISSING{ablation.model} (full model PR-AUC \MISSING{ablation.full}) changed PR-AUC by \MISSING{ablation.metadata} (metadata), \MISSING{ablation.developer} (developer), \MISSING{ablation.permissions} (permissions), \MISSING{ablation.listing\_text} (listing text) and \MISSING{ablation.review\_text} (review text) (Fig.~\ref{fig:ablation}). The largest drop came from removing \MISSING{ablation.most\_important}.

\begin{figure}[t]
\centering
\includegraphics[width=0.7\linewidth]{figures/ablation.png}
\caption{Change in PR-AUC on organic apps when each feature group is removed (error bars: fold standard deviation).}\label{fig:ablation}
\end{figure}

% ---------------------------------------------------------------------
\subsection{Early detection}\label{sec:res-early}

Using only information available in an app's first 14 days (no install counts, ratings or age, and reviews from the first two weeks only), \MISSING{early.model} reached a PR-AUC of \MISSING{early.early14.pr\_auc} (ROC-AUC \MISSING{early.early14.roc\_auc}), against \MISSING{early.full.pr\_auc} (ROC-AUC \MISSING{early.full.roc\_auc}) with full snapshot features. Only \MISSING{early.n\_with\_reviews} apps had any review in their first 14 days, because reviews were collected newest-first, so the early setting relies mainly on listing text and permissions.

\AUTHOR{State whether early detection is practical and which signals carry it.}

% ---------------------------------------------------------------------
\subsection{Class imbalance and leakage checks}\label{sec:res-checks}

\begin{table}[t]
\centering
\caption{PR-AUC on organic apps: class weights vs SMOTE (left) and random vs developer-grouped CV (right).}\label{tab:checks}
\begin{tabular}{lcc@{\hskip 18pt}ccc}
\toprule
Model & Weights & SMOTE & Grouped & Random & Gap \\
\midrule
Logistic regression & \MISSING{smote.logreg.weights} & \MISSING{smote.logreg.smote} & \MISSING{leakage.logreg.grouped} & \MISSING{leakage.logreg.random} & \MISSING{leakage.logreg.gap} \\
Random forest & \MISSING{smote.rf.weights} & \MISSING{smote.rf.smote} & \MISSING{leakage.rf.grouped} & \MISSING{leakage.rf.random} & \MISSING{leakage.rf.gap} \\
XGBoost & \MISSING{smote.xgb.weights} & \MISSING{smote.xgb.smote} & \MISSING{leakage.xgb.grouped} & \MISSING{leakage.xgb.random} & \MISSING{leakage.xgb.gap} \\
LightGBM & \MISSING{smote.lgbm.weights} & \MISSING{smote.lgbm.smote} & \MISSING{leakage.lgbm.grouped} & \MISSING{leakage.lgbm.random} & \MISSING{leakage.lgbm.gap} \\
\bottomrule
\end{tabular}
\end{table}

Table~\ref{tab:checks} compares class weighting with SMOTE oversampling (applied inside training folds only), and random with developer-grouped splits. A positive gap means random splitting overstates performance because near-identical apps from one developer leak across folds.

\AUTHOR{One sentence on each comparison: does SMOTE help, and how large is the leakage effect?}

% ---------------------------------------------------------------------
\subsection{Explanations}\label{sec:res-shap}

Figure~\ref{fig:shap} shows SHAP values for \MISSING{shap.model} fitted on all \MISSING{shap.n\_apps} labelled apps. The five most influential features were \MISSING{shap.top1}, \MISSING{shap.top2}, \MISSING{shap.top3}, \MISSING{shap.top4} and \MISSING{shap.top5}.

\begin{figure}[t]
\centering
\includegraphics[width=0.62\linewidth]{figures/shap_summary_beeswarm.png}
\caption{SHAP summary for \MISSING{shap.model}. Points to the right push the prediction toward \emph{unverified}; colour shows the feature value.}\label{fig:shap}
\end{figure}

\AUTHOR{Interpret the top features. Expected from EDA: unverified apps less often name an NBFC/bank partner or state the tenure, use free email addresses more often, and are newer; RBI-listed lenders request SMS access more often. Note that app age may partly reflect survivorship after the Play removal.}

For the review-text model, the n-grams most associated with unverified apps included \MISSING{text.top\_unverified}, and those most associated with verified apps included \MISSING{text.top\_verified} (Fig.~\ref{fig:text}).

\begin{figure}[t]
\centering
\includegraphics[width=0.62\linewidth]{figures/shap_text_top_terms.png}
\caption{Review n-grams with the largest logistic-regression weights toward each class.}\label{fig:text}
\end{figure}

\AUTHOR{Comment on whether these are risk terms or app/product names (the preliminary run showed app names such as ``kuhoo'', which means the text model partly memorises apps).}

% ---------------------------------------------------------------------
\subsection{Summary}\label{sec:res-summary}

\AUTHOR{Three or four sentences: best model and its organic PR-AUC versus the base rate; what the temporal and early-detection results say about flagging new apps; the most informative signals; the main caveat (labels mean ``not in the RBI list'', not ``illegal'').}
